\section{Introduction}
\label{sec:intro}

Immersive applications increasingly define the trajectory from \fiveg
to \gls{6g}~\cite{liberg2024,shen2023toward}: 3GPP and ITU-R list
immersive and extended-reality services among the primary usage
scenarios of these generations~\cite{3gpp22261,itur-m2160}, and mobile
video traffic is projected to more than double by
2031~\cite{mobility25nov}. The most demanding of these services is
immersive media streaming for \gls{vr}: volumetric content captured from
a real scene and reproduced in an \gls{hmd} with
\gls{6dof}~\cite{Immersivevideo,Viola2023}. Its dominant format, the
dynamic point cloud, requires gigabits per second
uncompressed~\cite{pcstreaming21}, and \gls{mpeg}'s
\gls{vpcc}~\cite{v3c,vpccgraziosi} is the reference codec for dense,
dynamic content. \vpcc codes geometry, attribute and occupancy atlases
with \gls{hevc}~\cite{hevc} and inherits its temporal prediction
structure: the bitstream is organised in \glspl{gof}, each of which must
be received in full before any of its frames can be decoded. Because
standalone \glspl{hmd} cannot decode \vpcc in real time, deployments
render at an edge server and cast the frames to the headset, so the
path from the content origin to the edge renderer is where network
impairment reaches the content.

Volumetric streaming research has largely assumed lossy datagram
transport or partially reliable \quic variants that drop stale
data~\cite{ravuri2023,haseeb2025}. Production systems run neither: the
\quic shipped in browsers, content delivery networks and mobile
operating systems is fully reliable~\cite{rfc9000}. No bitstream reaches
the decoder incomplete; loss becomes delay through retransmission and
congestion-window reduction, and the \gof amplifies that delay, since
one late packet stalls the whole unit. The perceptual consequence is
temporal rather than geometric: freezes and frame skips, not holes in
the surface. Our earlier work showed that latency and loss leave
distinct transport signatures~\cite{vquest2026}, and that a \vpcc slice
with \emph{lower} loss but a tighter \gof deadline can score lower
than one with higher loss~\cite{mosaic2026}. Existing \qoe models do not
capture this regime. The closest, by Udora et
al.~\cite{qoeholographic,udora2026tmm}, fits an exponential-in-delay,
linear-in-loss regression per quality level and reaches held-out
correlations of 0.57--0.66; it describes its data but not \emph{why}
delay and loss degrade \qoe, and so cannot transfer to a new \gof size,
encoder or network. Subjective \vpcc studies, meanwhile, inject loss
offline or assess outside fully occlusive
\gls{vr}~\cite{dumic2023,nguyen2023,Cox2023}.

We present \crystal, a mechanistic two-layer \qoe model for \vpcc
streaming over reliable \quic. Layer~1, which has no perceptual
parameter, derives the \emph{temporal loss fraction} $\tlf$, the fraction
of wall-clock time in which the renderer has no fresh \gof to show, from
delivery slack, round-trip time, loss-recovery cost and \gof size.
Layer~2 maps $\tlf$ to opinion score with one constant,
\begin{equation}
  \crystalmodel,
  \label{eq:crystal-intro}
\end{equation}
where $\qc$ is the clean-reference ceiling of each content and $\alpha$
is the only fitted parameter. The model claims that, under reliable
transport, latency and loss act on \qoe only through $\tlf$. We test it
with an ITU-T P.910 \gls{acr} study ($N=50$) on a private \fiveg testbed
that streams \vpcc over \gls{http3} to an edge renderer and casts the
rendered frames to Meta Quest~3 headsets, using a $2\times2$ factorial
design that separates latency, loss and their interaction at a fixed
encoder configuration. The contributions are:
\begin{enumerate}[leftmargin=*,nosep]
  \item \textbf{A mechanistic model} whose transport layer is
        parameter-free and whose perceptual layer has one constant,
        replacing per-quality-level regression coefficients with a
        single explanatory variable, $\tlf$.
  \item \textbf{A pre-registered factorial evaluation} showing that
        latency and loss each degrade perceived smoothness and together
        degrade it $1.61\times$ more than the sum of the two, refuting
        the additivity of prior regressions; that \gof structure rather
        than loss magnitude governs resilience, a larger low-delay \gof
        being more robust than a smaller one; and that a single
        $\alpha$ \todo{does / does not} predict held-out conditions.
  \item \textbf{A transport characterisation and dataset}: per-\gof
        stall and deadline-miss measurements of reliable \quic under
        controlled impairment, and \todo{$N$} ratings over three
        \gls{mpeg} test sequences collected in \gls{vr} with a
        reproducible rating application.
\end{enumerate}

\blind{%
This article extends two conference papers: \cite{mosaic2026}
introduced the \vpcc-over-\quic pipeline and a three-tier study
($N=20$) that paired impairment with encoder configuration, and
\cite{vquest2026} characterised the transport signatures of latency and
loss for a frame-independent codec; earlier testbed descriptions
appeared as demonstrations~\cite{bestdemo25,nfvsdn24}. Neither proposed
a model or isolated \gof structure.%
}{%
This article extends two conference papers by the authors:
\mosaic~\cite{mosaic2026} introduced the \vpcc-over-\quic pipeline and
a three-tier study ($N=20$) that paired impairment with encoder
configuration, and \vquest~\cite{vquest2026} characterised the
transport signatures of latency and loss for a frame-independent codec;
earlier testbed descriptions appeared as
demonstrations~\cite{bestdemo25,nfvsdn24}. Neither proposed a model or
isolated \gof structure.%
}
Section~\ref{sec:related} reviews related work,
Section~\ref{sec:model} derives the model and its predictions,
Section~\ref{sec:system} describes the testbed and experimental design,
Section~\ref{sec:results} reports the results,
Section~\ref{sec:discussion} discusses implications and limitations,
and Section~\ref{sec:conclusion} concludes.
